% Source: Wald GR, physical PDF pp. 463--471
%         (printed pp. 450--458).
% Coverage: Appendix E.1, Lagrangian Formulation; equations (E.1.1)--(E.1.42).
% Figures/tables/footnotes: no figures/tables; footnote 1.
% Status: source-reviewed and compile-clean.
% Uncertainties: none.

\subsection{Lagrangian Formulation}\label{sec:E.1}

The dynamical content of general relativity is fully expressed by Einstein's field
equation, \(G_{ab}=8\pi T_{ab}\). Nevertheless, even in a purely classical
(i.e., non-quantum) context, it is convenient and useful for many purposes to have
Lagrangian and Hamiltonian formulations of general relativity. For example, as we
shall see below, Einstein's equation can be derived from a very simple and natural
Lagrangian, thus contributing further to the aesthetic appeal of general relativity.
The Hamiltonian formulation yields insights into the nature of dynamics of general
relativity. Indeed, although we analyzed the initial value formulation of general
relativity in chapter 10 using only the field equation, the viewpoint that Einstein's
equation describes the evolution of the spatial metric, \(h_{ab}\), with "time" is
perhaps best motivated via the Hamiltonian formulation. In addition, the Hamiltonian
formulation also motivates the definition given in chapter 11 of the total energy at
spatial infinity of an asymptotically flat spacetime.

However, an even stronger reason for studying the Lagrangian and Hamiltonian
formulations of general relativity arises from the desire to obtain a quantum theory
of gravitation. Although the entire content of a classical field theory is expressed
by the field equation, most prescriptions for formulating a quantum field theory
associated with the classical theory require that the classical theory be expressed
in a Lagrangian or Hamiltonian form. In particular, the path integral formulation
requires that one have an action principle at the classical level---i.e., it requires
a Lagrangian formulation of the classical theory---while the canonical quantization
procedure requires that the classical field theory be cast in Hamiltonian form. Thus,
the Lagrangian and Hamiltonian formulations of general relativity may play an
important role in the development of a quantum theory of gravity (see chapter 14).

We begin our discussion by explaining what we mean by a Lagrangian formulation
of a field theory. Consider a theory involving a tensor field (or collection of tensor
fields) defined on a manifold \(M\). We shall suppress all indices and denote the
field (or fields) by \(\psi\). Let \(S[\psi]\) be a functional of \(\psi\), i.e., \(S\)
is a map from field configurations on \(M\) into numbers. Let \(\psi_\lambda\) be a
smooth one parameter family of field configurations starting from \(\psi_0\) which
satisfy appropriate boundary conditions. We denote
\(\left.\dd\psi_\lambda/\dd\lambda\right|_{\lambda=0}\) by \(\delta\psi\). Suppose
\(\dd S/\dd\lambda\) at \(\lambda=0\) exists for all such one-parameter families
starting from \(\psi_0\). Suppose, furthermore, that there exists a smooth tensor
field \(\chi\) [which is dual to \(\psi\), i.e., if \(\psi\) is a tensor field of type
\((k,l)\), then \(\chi\) will be of type \((l,k)\)] such that for all such families
we have
\begin{equation}
  \left.\frac{\dd S}{\dd\lambda}\right|_{\lambda=0}
    =\int_M \chi\,\delta\psi .
  \tag{E.1.1}\label{eq:E.1.1}
\end{equation}
where contraction of all indices in the integral is understood. Then we say that
\(S\) is \emph{functionally differentiable} at \(\psi_0\). We call \(\chi\) the
\emph{functional derivative}\footnote{More generally, if there exists a tensor
distribution \(\chi\) such that
\(\left.\dd S/\dd\lambda\right|_{\lambda=0}=\chi[\delta\psi]\), we also say that \(S\)
is functionally differentiable and refer to \(\chi\) as the functional derivative of
\(S\) at \(\psi_0\).} of \(S\) and denote it as
\begin{equation}
  \chi=\left.\frac{\delta S}{\delta\psi}\right|_{\psi_0}.
  \tag{E.1.2}\label{eq:E.1.2}
\end{equation}

Consider, now, a functional \(S\) of the form
\begin{equation}
  S[\psi]=\int_M\mathcal L[\psi],
  \tag{E.1.3}\label{eq:E.1.3}
\end{equation}
where \(\mathcal L\) is a local function of \(\psi\) and a finite number of its
derivatives, i.e.,
\begin{equation}
  \left.\mathcal L\right|_x
    =\mathcal L\bigl(\psi(x),\nabla\psi(x),\ldots,\nabla^k\psi(x)\bigr).
  \tag{E.1.4}\label{eq:E.1.4}
\end{equation}
Suppose that \(S\) is functionally differentiable and that the field configurations
\(\psi\) which extremize \(S\),
\begin{equation}
  \left.\frac{\delta S}{\delta\psi}\right|_\psi=0,
  \tag{E.1.5}\label{eq:E.1.5}
\end{equation}
are precisely the ones which are solutions of the field equation for \(\psi\). Then
\(S\) is called an \emph{action}, \(\mathcal L\) is called a \emph{Lagrangian
density}, and the specification of such an \(\mathcal L\) is what we mean by a
Lagrangian formulation of the field theory.

Thus the notion of a Lagrangian formulation of a field theory is closely analogous
to that of a Lagrangian formulation in ordinary particle mechanics. In particle
mechanics one specifies an action functional of the particle path as an integral of a
Lagrangian function over the path. The variational problem analogous to
Eq.~\eqref{eq:E.1.5} is made precise by focusing attention on paths of finite length
% SOURCE ERRATUM: printed ``wih''; corrected to ``with.''
and extremizing the action with respect to path variations which keep the path
endpoints fixed. By analogy, to make our variational problem in the field case
precise, we shall focus attention on a compact region, \(U\), of \(M\) and will
consider one-parameter families \(\psi_\lambda\) which keep fixed the value of
\(\psi\) on the boundary, \(\partial U\).

As a simple example, we give a Lagrangian formulation of the theory of a
Klein--Gordon scalar field \(\phi\) in Minkowski spacetime. We define
\begin{equation}
  \mathcal L_{\mathrm{KG}}
    =-\frac{1}{2}\bigl(\partial_a\phi\partial^a\phi+m^2\phi^2\bigr).
  \tag{E.1.6}\label{eq:E.1.6}
\end{equation}
(Here, the normalization of \(\mathcal L_{\mathrm{KG}}\) is chosen for later
convenience in defining conjugate momenta [see section E.2].) We obtain
\begin{align}
  \left.\frac{\dd S_{\mathrm{KG}}}{\dd\lambda}\right|_{\lambda=0}
    &= -\int\bigl[\partial_a\phi_0\partial^a(\delta\phi)
        +m^2\phi_0(\delta\phi)\bigr] \notag\\
    &= \int\bigl[\partial_a\partial^a\phi_0-m^2\phi_0\bigr]\delta\phi .
  \tag{E.1.7}\label{eq:E.1.7}
\end{align}
where the natural volume element on Minkowski spacetime is understood in the
integrals, and the boundary term in the integration by parts does not contribute on
account of our boundary condition on \(\phi_\lambda\), which requires
\(\delta\phi=0\) on the boundary. Thus, \(S_{\mathrm{KG}}\) is functionally
differentiable and we have
\begin{equation}
  \frac{\delta S_{\mathrm{KG}}}{\delta\phi}
    =\partial_a\partial^a\phi-m^2\phi.
  \tag{E.1.8}\label{eq:E.1.8}
\end{equation}
Consequently, Eq.~\eqref{eq:E.1.5} is just the Klein--Gordon equation
\eqref{eq:4.2.19}, as desired. Similarly, the function
\begin{equation}
  \mathcal L_{\mathrm{EM}}
    =-\frac{1}{4}F_{ab}F^{ab}
    =-\partial_{[a}A_{b]}\partial^{[a}A^{b]}
  \tag{E.1.9}\label{eq:E.1.9}
\end{equation}
of the field variable \(A_a\) is a Lagrangian density for Maxwell's equations in
Minkowski spacetime. Indeed, Eqs.~\eqref{eq:E.1.6} and \eqref{eq:E.1.9} often
are taken as the starting point for the study of Klein--Gordon and Maxwell fields:
One writes down the simple and natural scalar functions \(\mathcal L_{\mathrm{KG}}\)
or \(\mathcal L_{\mathrm{EM}}\) and obtains the field equations via
Eq.~\eqref{eq:E.1.5}.

For general relativity, the field variable is the spacetime metric, \(g_{ab}\),
defined on a four-dimensional manifold, \(M\). In this case, a slight awkwardness
results from the fact the natural volume element to use in the integrals
\eqref{eq:E.1.1} and \eqref{eq:E.1.3} is the volume element \(\epsilon_{abcd}\)
determined from \(g_{ab}\) via Eq.~\eqref{eq:B.2.9}. Consequently, the volume
element itself depends on the field variable, and hence its variation must be taken
into account when calculating functional derivatives. One way to handle this
situation would be to define \(\mathcal L\) to be a totally antisymmetric four-index
tensor rather than a scalar, i.e., to incorporate the volume element into
\(\mathcal L\). This would require us to make a similar modification of our
definition of functional derivatives. Instead we shall follow the considerably less
cumbersome procedure of introducing a fixed volume element
\(e_{abcd}=e_{[abcd]}\) on \(M\) and defining all integrals over \(M\) to be with
respect to \(e_{abcd}\) rather than \(\epsilon_{abcd}\). One way to do this (at
least over a portion of \(M\)) would be to choose a coordinate system and take
\(e_{abcd}\) to be the associated coordinate volume element, but we emphasize that
the introduction of a coordinate system is not necessary. Since any two volume
elements differ at each point by at most a scalar factor, we have
\begin{equation}
  \epsilon_{abcd}=f e_{abcd}.
  \tag{E.1.10}\label{eq:E.1.10}
\end{equation}
In any basis where the nonvanishing components of \(e_{abcd}\) have the values
\(\pm1\), the calculation which led to Eq.~\eqref{eq:B.2.16} for the case of a
coordinate basis and its associated volume element shows that \(f=\sqrt{-g}\),
where \(g\) denotes the determinant of the matrix of components, \(g_{\mu\nu}\), of
the metric in that basis. Hence, we shall follow the same notational convention as
used for coordinate bases in section 3.4a and denote \(f\) as \(\sqrt{-g}\). Given
the volume element \(e_{abcd}\) on \(M\), we define a \emph{tensor density}
\(T^{a\ldots b}{}_{c\ldots \dd}\) to be a tensor which can be expressed in the form
\begin{equation}
  T^{a\ldots b}{}_{c\ldots \dd}
    =\sqrt{-g}\,\widetilde T^{a\ldots b}{}_{c\ldots \dd},
  \tag{E.1.11}\label{eq:E.1.11}
\end{equation}
where \(\widetilde T^{a\ldots b}{}_{c\ldots \dd}\) is a tensor whose value does not
depend on the choice of \(e_{abcd}\). In order that the action, \(S\), of general
relativity be independent of \(e_{abcd}\), it is necessary that the Lagrangian
density \(\mathcal L\) be a scalar density. Similarly, in order for \(\dd S/\dd\lambda\)
to be independent of \(e_{abcd}\), the functional derivatives of \(S\) must be
tensor densities.

We now shall demonstrate that---except for boundary terms to be dealt with
later---the scalar density
\begin{equation}
  \mathcal L_G=\sqrt{-g}\,R
  \tag{E.1.12}\label{eq:E.1.12}
\end{equation}
is a Lagrangian density for the vacuum Einstein equation. The corresponding action
\begin{equation}
  S[g^{ab}]=\int\mathcal L_G e
  \tag{E.1.13}\label{eq:E.1.13}
\end{equation}
is known as the \emph{Hilbert action}. Here we have written \(e\) (using the
differential forms notation of appendix B) in order to emphasize our use of this
volume element. In addition, for convenience, we have taken the inverse metric
\(g^{ab}\) as the field variable rather than \(g_{ab}\). For a one-parameter
variation we define \(\delta g^{ab}\) to be \(\dd g^{ab}/\dd\lambda\). However, in
order to use without modification the results of section 7.5 where \(g_{ab}\) was
used as the independent variable, we shall define
\(\delta g_{ab}=\dd g_{ab}/\dd\lambda\). Thus, we warn the reader that, since
\(g^{ac}g_{cb}=\delta^a{}_b\), we have
\(\delta g_{ab}=-g_{ac}g_{bd}\delta g^{cd}\), i.e., in this section we will
\emph{not} use the unperturbed metric to raise and lower indices of the metric
perturbations. Note also that since \(g^{ab}\) and hence \(\delta g^{ab}\) must be
symmetric, one can add an antisymmetric tensor to any functional derivative with
respect to \(g^{ab}\) without affecting Eq.~\eqref{eq:E.1.1}. We eliminate this
arbitrariness by requiring that all such functional derivatives be symmetric.

For a one-parameter family starting from \(g^{ab}\), we have
\begin{equation}
  \frac{\dd\mathcal L_G}{\dd\lambda}
    =\sqrt{-g}(\delta R_{ab})g^{ab}
      +\sqrt{-g}R_{ab}\delta g^{ab}+R\delta(\sqrt{-g}).
  \tag{E.1.14}\label{eq:E.1.14}
\end{equation}
From Eq.~\eqref{eq:7.5.14} of chapter 7, we have
\begin{equation}
  g^{ab}\delta R_{ab}=\nabla^a v_a,
  \tag{E.1.15}\label{eq:E.1.15}
\end{equation}
where
\begin{equation}
  v_a=\nabla^b(\delta g_{ab})-g^{cd}\nabla_a(\delta g_{cd}).
  \tag{E.1.16}\label{eq:E.1.16}
\end{equation}
In addition, using Eq.~\eqref{eq:9.3.11}, we have
\begin{align}
  \delta(\sqrt{-g})
    &=\frac{1}{2}\sqrt{-g}\,g^{ab}\delta g_{ab}\notag\\
    &=-\frac{1}{2}\sqrt{-g}\,g_{ab}\delta g^{ab}.
  \tag{E.1.17}\label{eq:E.1.17}
\end{align}
Thus we obtain
\begin{equation}
  \frac{\dd S_G}{\dd\lambda}
    =\int\frac{\dd\mathcal L_G}{\dd\lambda}e
    =\int\nabla^a v_a\sqrt{-g}\,e
      +\int\left(R_{ab}-\frac{1}{2}Rg_{ab}\right)
        \delta g^{ab}\sqrt{-g}\,e.
  \tag{E.1.18}\label{eq:E.1.18}
\end{equation}
The first term in Eq.~\eqref{eq:E.1.18} is the integral of a divergence,
\(\nabla^a v_a\), with respect to the natural volume element
\(\epsilon=\sqrt{-g}\,e\). Hence, by Stokes's theorem this integral contributes
only a boundary term. In fact, this term does not vanish for general variations where
\(g^{ab}\) is held fixed on the boundary, although it does vanish for variations where
the first derivatives of \(g^{ab}\) also are held fixed. However, in order to simplify
the discussion here, we shall ignore this contribution for the present. (At the end of
this section we shall calculate this term and show how to modify \(S_G\) to cancel
its contribution.) Thus discarding the boundary term, we find
\begin{equation}
  \frac{\delta S_G}{\delta g^{ab}}
    =\sqrt{-g}\left(R_{ab}-\frac{1}{2}Rg_{ab}\right),
  \tag{E.1.19}\label{eq:E.1.19}
\end{equation}
and Eq.~\eqref{eq:E.1.5} is seen to be equivalent to Einstein's equation in vacuum,
as desired. Thus, from the Lagrangian viewpoint, Einstein's equation arises in a
very natural way, since the Lagrangian density of Eq.~\eqref{eq:E.1.12} is
unquestionably one of the simplest scalar densities which can be constructed from
the spacetime metric.

It is interesting to note that instead of viewing the metric alone as the field variable
for general relativity, we could view the metric and the derivative operator
\(\nabla_a\) as independent variables. Remarkably, if we use the same Lagrangian
density \eqref{eq:E.1.12} but now view \(R_{ab}\) as a function of the derivative
operator alone (i.e., independent of \(g^{ab}\)) and vary the \emph{Palatini action},
\begin{equation}
  \mathcal S_G[g^{ab},\nabla_a]
    =\int\sqrt{-g}R_{ab}g^{ab}e,
  \tag{E.1.20}\label{eq:E.1.20}
\end{equation}
with respect to both \(g^{ab}\) and \(\nabla_a\), we recover Einstein's equation
\eqref{eq:E.1.19} together with the metric compatibility condition
\(\nabla_cg^{ab}=0\) on the derivative operator. To prove this we begin by noting
that since \(\nabla_a\) can be expressed in terms of an arbitrary fixed derivative
operator \(\widetilde\nabla_a\) and a tensor field \(C^c{}_{ab}\) (see section 3.1),
variation of \(\nabla_a\) is equivalent to variation of \(C^c{}_{ab}\). In considering
one parameter variations of \(g^{ab}\) and \(\nabla_a\), it will be convenient to
choose \(\widetilde\nabla_a\) to be the derivative operator compatible with
\(g^{ab}\) at \(\lambda=0\). The key change from the previous calculation is that
we must use Eq.~\eqref{eq:7.5.10} rather than Eq.~\eqref{eq:7.5.14} to evaluate
\(\delta R_{ab}\). We find at \(\lambda=0\)
\begin{align}
  \frac{\dd\mathcal S_G}{\dd\lambda}
    &=-2\int g^{ab}\nabla_{[a}\delta C^c{}_{c]b}\sqrt{-g}\,e
      +\int\left(R_{ab}-\frac{1}{2}Rg_{ab}\right)
        \delta g^{ab}\sqrt{-g}\,e\notag\\
    &=-2\int g^{ab}\widetilde\nabla_{[a}\delta C^c{}_{c]b}\sqrt{-g}\,e
      +\int g^{ab}\bigl[
        C^d{}_{ab}\delta C^c{}_{cd}
        +C^c{}_{cd}\delta C^d{}_{ab}
        -2C^d{}_{cb}\delta C^c{}_{ad}\bigr]\sqrt{-g}\,e\notag\\
    &\quad+\int\left(R_{ab}-\frac{1}{2}Rg_{ab}\right)
        \delta g^{ab}\sqrt{-g}\,e\notag\\
    &=\int\bigl[
        C^b{}_{\dd d}\delta^a{}_c
        +C^d{}_{\dd c}g^{ab}
        -2C^b{}_{ca}\bigr]\delta C^c{}_{ab}\sqrt{-g}\,e
      +\int\left(R_{ab}-\frac{1}{2}Rg_{ab}\right)
        \delta g^{ab}\sqrt{-g}\,e .
  \tag{E.1.21}\label{eq:E.1.21}
\end{align}
Here, the first term on the right side of the second line vanishes by Stokes's
theorem, since \(\widetilde\nabla_a\) is the metric compatible derivative operator.
(In this case there is no boundary term since we require \(\delta C^c{}_{ab}\) to
vanish on the boundary.) The vanishing of
\(\delta\mathcal S_G/\delta C^c{}_{ab}\) requires the term in brackets in the final
equality of Eq.~\eqref{eq:E.1.21} to vanish when symmetrized over \(a\) and \(b\),
which, after some algebra, implies \(C^c{}_{ab}=0\), i.e.,
\(\nabla_a=\widetilde\nabla_a\). The vanishing of
\(\delta\mathcal S_G/\delta g^{ab}\) yields Einstein's equation as before.

The non-vacuum Einstein equation with matter fields such as the Klein--Gordon
scalar field or the Maxwell field also can be obtained from a Lagrangian formulation
in a very simple and natural manner. First, we must find a suitable Lagrangian
density \(\mathcal L_M\) for the matter fields in curved spacetime. In particular,
for the Klein--Gordon field it is easily verified that functional differentiation with
respect to \(\phi\) of the action, \(S_{\mathrm{KG}}\), obtained from the Lagrangian
density,
\begin{equation}
  \mathcal L_{\mathrm{KG}}
    =-\frac{1}{2}\sqrt{-g}
      \bigl(g^{ab}\nabla_a\phi\nabla_b\phi+m^2\phi^2\bigr),
  \tag{E.1.22}\label{eq:E.1.22}
\end{equation}
yields the Klein--Gordon equation \eqref{eq:4.3.9} in curved spacetime. Similarly,
\begin{equation}
  \mathcal L_{\mathrm{EM}}
    =-\frac{1}{4}\sqrt{-g}\,g^{ac}g^{bd}F_{ab}F_{cd}
    =-\sqrt{-g}\,g^{ac}g^{bd}\nabla_{[a}A_{b]}\nabla_{[c}A_{d]}
  \tag{E.1.23}\label{eq:E.1.23}
\end{equation}
yields Maxwell's equations in curved spacetime. To obtain the coupled
Einstein--matter field equations, we construct a (total) Lagrangian density,
\(\mathcal L\), by adding together the Einstein Lagrangian density \(\mathcal L_G\)
with a multiple of the Lagrangian density, \(\mathcal L_M\), for the matter field,
\begin{equation}
  \mathcal L=\mathcal L_G+\alpha_M\mathcal L_M,
  \tag{E.1.24}\label{eq:E.1.24}
\end{equation}
where \(\alpha_M\) is a constant. Since \(\mathcal L_G\) does not depend on the
matter field, variation of the total action, \(S\), with respect to it will yield the
same equation as variation of \(S_M\) alone. Variation of \(S\) with respect to
\(g^{ab}\) yields the equation
\begin{equation}
  G_{ab}=R_{ab}-\frac{1}{2}Rg_{ab}=8\pi T_{ab},
  \tag{E.1.25}\label{eq:E.1.25}
\end{equation}
where the tensor \(T_{ab}\) is given by
\begin{equation}
  T_{ab}=-\frac{\alpha_M}{8\pi}\frac{1}{\sqrt{-g}}
    \frac{\delta S_M}{\delta g^{ab}}.
  \tag{E.1.26}\label{eq:E.1.26}
\end{equation}
For the Lagrangian densities \eqref{eq:E.1.22} and \eqref{eq:E.1.23}, it is easily
verified that for appropriate choices of \(\alpha_M\), \(T_{ab}\) agrees, respectively,
with Eqs.~\eqref{eq:4.3.10} and \eqref{eq:4.3.14}. Thus, the Lagrangian density
\eqref{eq:E.1.24} with \(\mathcal L_M=\mathcal L_{\mathrm{KG}}\) and
\(\alpha_{\mathrm{KG}}=16\pi\) yields the coupled Einstein--Klein--Gordon
equations, whereas \eqref{eq:E.1.24} with
\(\mathcal L_M=\mathcal L_{\mathrm{EM}}\) and \(\alpha_{\mathrm{EM}}=4\) yields the
Einstein--Maxwell equations. More generally, if one takes a Lagrangian density
\(\mathcal L_M\) as the starting point in the definition of a matter field theory,
then Eq.~\eqref{eq:E.1.26} may be used to \emph{define} the stress-energy tensor,
\(T_{ab}\), of that field. If the Lagrangian density for the matter field does not
depend on the choice of derivative operator \(\nabla_a\), then the Einstein--matter
equations also may be derived from variation of the sum of the Palatini action and
matter action.

The matter action \(S_M\) must be invariant under diffeomorphisms, i.e., if
\(f_\lambda:M\to M\) is a one-parameter family of diffeomorphisms, we have
\(S_M[g^{ab},\psi]=S_M[f_\lambda^*g^{ab},f_\lambda^*\psi]\), where
\(f_\lambda^*\) is defined in appendix C. Hence, for such variations, we have
\begin{equation}
  0=\frac{\dd S_M}{\dd\lambda}
    =\int\frac{\delta S_M}{\delta g^{ab}}\delta g^{ab}
      +\int\frac{\delta S_M}{\delta\psi}\delta\psi .
  \tag{E.1.27}\label{eq:E.1.27}
\end{equation}
Recall from appendix C that for such variations, \(\delta g^{ab}\) has the general
form \(\mathcal L_w g^{ab}=2\nabla^{(a}w^{b)}\), where \(w^a\) is an arbitrary
vector field. Suppose, now, that \(\psi\) satisfies the matter field equations. Then
\(\left.\delta S_M/\delta\psi\right|_\psi=0\) and the second term in
Eq.~\eqref{eq:E.1.27} makes no contribution. Thus, using the definition
\eqref{eq:E.1.26}, we find that if \(\psi\) satisfies the matter field equations,
then for all smooth \(w^a\) of compact support, we have
\begin{align}
  0
    &=\int\sqrt{-g}T_{ab}\nabla^{(a}w^{b)}e\notag\\
    &=\int T_{ab}\nabla^a w^b\epsilon\notag\\
    &=-\int(\nabla^aT_{ab})w^b\epsilon ,
  \tag{E.1.28}\label{eq:E.1.28}
\end{align}
which implies that
\begin{equation}
  \nabla^aT_{ab}=0.
  \tag{E.1.29}\label{eq:E.1.29}
\end{equation}
Thus, for a diffeomorphism invariant action, \(T_{ab}\) always is conserved by
virtue of the matter field equation. This reinforces the interpretation of \(T_{ab}\)
as representing the stress-energy-momentum tensor of the matter field. Note also
that by applying the above argument to \(S_G\), it follows that (independent of any
field equation) we have
\begin{equation}
  \nabla^aG_{ab}=0.
  \tag{E.1.30}\label{eq:E.1.30}
\end{equation}
Thus, in the Lagrangian formulation of general relativity, the contracted Bianchi
identity may be viewed as a consequence of the invariance of the Hilbert action
under diffeomorphisms.

In Minkowski spacetime \((\mathbb R^4,\eta_{ab})\) there exists an alternative
procedure for defining a stress-energy tensor associated with a field \(\psi\)
starting from its Lagrangian \(\mathcal L\). Consider, for simplicity, the case where
\(\mathcal L\) is a local function of \(\eta_{ab}\), \(\psi\), and
\(\partial_a\psi\) but no higher derivatives of \(\psi\). Then, the equation of
motion for \(\psi\) is
\begin{equation}
  0=\frac{\delta S}{\delta\psi}
    =\frac{\partial\mathcal L}{\partial\psi}
      -\partial_a\left[\frac{\partial\mathcal L}
        {\partial(\partial_a\psi)}\right].
  \tag{E.1.31}\label{eq:E.1.31}
\end{equation}
Here, in the case where \(\psi\) is a scalar field
\(\partial\mathcal L/\partial(\partial_a\psi)\) means the vector field \(V^a\)
which at each point \(x\) satisfies
\begin{equation}
  \frac{\dd\mathcal L}{\dd\lambda}=V^a\delta(\partial_a\psi)
  \tag{E.1.32}\label{eq:E.1.32}
\end{equation}
for all smooth one-parameter variations of \(\psi\) which keep fixed at \(x\) the
value of all quantities upon which \(\mathcal L\) depends except
\(\partial_a\psi\). (The components of \(V^a\) in a basis are just the partial
derivatives of \(\mathcal L\) with respect to the corresponding dual basis
components of \(\partial_a\psi\).) If \(\psi\) is a tensor of type \((k,l)\), then
\(V^a=\partial\mathcal L/\partial(\partial_a\psi)\) will be a tensor of type
\((l+1,k)\), with contraction of all indices understood in Eq.~\eqref{eq:E.1.32}.
Partial derivatives of \(\mathcal L\) with respect to other tensor variables are
defined similarly. For an arbitrary smooth one-parameter family
\((\psi_\lambda,(\eta_\lambda)_{ab})\) of fields \(\psi_\lambda\) and flat metrics
\((\eta_\lambda)_{ab}\), we have
\begin{equation}
  \delta\mathcal L\equiv\frac{\dd\mathcal L}{\dd\lambda}
    =\frac{\partial\mathcal L}{\partial\psi}\delta\psi
      +\frac{\partial\mathcal L}{\partial(\partial_a\psi)}
        \delta(\partial_a\psi)
      +\frac{\partial\mathcal L}{\partial\eta_{ab}}\delta\eta_{ab},
  \tag{E.1.33}\label{eq:E.1.33}
\end{equation}
where no boundary conditions on the varied quantities need be assumed here.
Consider, now, the variations of \(\mathcal L\) produced by a one-parameter family
of diffeomorphisms generated by a vector field \(\xi^a\). Then
Eq.~\eqref{eq:E.1.33} becomes
\begin{equation}
  \delta\mathcal L
    =\mathcal L_\xi\mathcal L
    =\xi^a\partial_a\mathcal L
    =\frac{\partial\mathcal L}{\partial\psi}\mathcal L_\xi\psi
      +\frac{\partial\mathcal L}{\partial(\partial_a\psi)}
        \mathcal L_\xi\partial_a\psi
      +\frac{\partial\mathcal L}{\partial\eta_{ab}}
        \mathcal L_\xi\eta_{ab}.
  \tag{E.1.34}\label{eq:E.1.34}
\end{equation}
Now restrict attention to the case where \(\xi^a\) is a Killing field. Then the
last term in Eq.~\eqref{eq:E.1.34} vanishes, and in the second term we have
\(\mathcal L_\xi\partial_a\psi=\partial_a(\mathcal L_\xi\psi)\). Using
Eq.~\eqref{eq:E.1.31} to substitute for
\(\partial\mathcal L/\partial\psi\), we obtain
\begin{equation}
  \partial_a\left[
    \frac{\partial\mathcal L}{\partial(\partial_a\psi)}\mathcal L_\xi\psi
    -\xi^a\mathcal L\right]=0.
  \tag{E.1.35}\label{eq:E.1.35}
\end{equation}
This result is known as \emph{Noether's theorem} as applied to the Poincaré group
of symmetries. In particular, the validity of Eq.~\eqref{eq:E.1.35} for all
translational Killing fields implies that the tensor
\begin{equation}
  S^{ab}=\frac{\partial\mathcal L}{\partial(\partial_a\psi)}
    \partial^b\psi-\eta^{ab}\mathcal L,
  \tag{E.1.36}\label{eq:E.1.36}
\end{equation}
known as the \emph{canonical energy-momentum tensor}, is conserved,
\(\partial_aS^{ab}=0\). For the Klein--Gordon field, we find that \(S_{ab}\) agrees
with \(T_{ab}\) (up to a numerical factor), where \(T_{ab}\) is defined by
Eq.~\eqref{eq:E.1.26} evaluated at Minkowski spacetime using the curved space
Lagrangian density \eqref{eq:E.1.22}. However, this agreement does not occur for
higher spin fields. Indeed, in the case of a Maxwell field, \(S^{ab}\) is not even
gauge invariant. Furthermore, in general \(S_{ab}\) is not symmetric, nor does it
naturally generalize to a conserved tensor in curved spacetime (Kuchař 1976). Thus
we adopt \(T_{ab}\) as our definition of the stress-energy tensor. It is the quantity
which naturally appears on the right-hand side of Einstein's equation
\eqref{eq:E.1.25} in a Lagrangian formulation of the Einstein--matter field
equations.

We conclude this section by evaluating the boundary term occurring in the variation
\eqref{eq:E.1.18} of the Hilbert action when \(\delta g^{ab}\) is required to vanish
on the boundary but no conditions are placed on the derivatives of \(\delta g^{ab}\).
We have
\begin{equation}
  \int_U\nabla_a v^a\epsilon=\int_{\partial U}v_an^a,
  \tag{E.1.37}\label{eq:E.1.37}
\end{equation}
where \(n^a\) is the unit normal to the boundary \(\partial U\) (which is assumed to
be non-null) and the natural volume element on \(\partial U\) is understood (see
appendix B). Using Eq.~\eqref{eq:E.1.16}, we have on \(\partial U\)
\begin{align}
  v_an^a
    &=n^ag^{bc}\bigl[\nabla_c(\delta g_{ab})
      -\nabla_a(\delta g_{bc})\bigr]\notag\\
    &=n^ah^{bc}\bigl[\nabla_c(\delta g_{ab})
      -\nabla_a(\delta g_{bc})\bigr]\notag\\
    &=-n^ah^{bc}\nabla_a(\delta g_{bc}),
  \tag{E.1.38}\label{eq:E.1.38}
\end{align}
where \(h_{ab}=g_{ab}\pm n_an_b\) is the induced metric on \(\partial U\) (see
chapter 10) and we have \(h^{bc}\nabla_c(\delta g_{ab})=0\) because
\(\delta g_{ab}=0\) on \(\partial U\). However, the right-hand side of
Eq.~\eqref{eq:E.1.38} is related to the variation of the trace of the extrinsic
curvature of the boundary. We have
\begin{equation}
  K\equiv K^a{}_a=h^a{}_b\nabla_a n^b
  \tag{E.1.39}\label{eq:E.1.39}
\end{equation}
and hence
\begin{align}
  \delta K
    &=h^a{}_b(\delta C)^b{}_{ac}n^c\notag\\
    &=\frac{1}{2}n^ch^a{}_bg^{bd}
      \bigl[\nabla_a(\delta g_{cd})
        +\nabla_c(\delta g_{ad})
        -\nabla_d(\delta g_{ac})\bigr]\notag\\
    &=\frac{1}{2}n^ch^{ad}\nabla_c(\delta g_{ad}).
  \tag{E.1.40}\label{eq:E.1.40}
\end{align}
Thus, under variations of the metric for which \(\delta g_{ab}=0\) on
\(\partial U\) we obtain from Eqs.~\eqref{eq:E.1.18}, \eqref{eq:E.1.38}, and
\eqref{eq:E.1.40}
\begin{equation}
  \frac{\dd S_G}{\dd\lambda}
    =-2\int_{\partial U}\delta K+\int_U G_{ab}\delta g^{ab}\epsilon .
  \tag{E.1.41}\label{eq:E.1.41}
\end{equation}
In fact, Eq.~\eqref{eq:E.1.41} continues to hold if we allow variations of \(g_{ab}\)
for which only the induced metric on the boundary is held fixed,
\(\delta h_{ab}=0\). This can be verified directly or deduced from the fact that if
\(\delta h_{ab}=0\) on the boundary, we can find a gauge transformation
\(\nabla_{(a}l_{b)}\) with \(l_b=0\) on the boundary which makes
\(\delta g_{ab}=0\). Since Eq.~\eqref{eq:E.1.41} holds for all variations with
\(\delta g_{ab}=0\) on \(\partial U\) and since all terms in Eq.~\eqref{eq:E.1.41}
are invariant under such gauge transformations, this equation must continue to hold
for variations which merely satisfy \(\delta h_{ab}=0\). Thus, the extremization of
\(S_G\) with respect to variations with \(\delta g_{ab}=0\) or \(\delta h_{ab}=0\)
on the boundary contains an additional, unwanted term. However, this can be
remedied by modifying \(S_G\). We define
\begin{equation}
  S'_G=S_G+2\int_{\partial U}K.
  \tag{E.1.42}\label{eq:E.1.42}
\end{equation}
Then extremization of \(S'_G\) yields the desired result since variation of the
boundary term in Eq.~\eqref{eq:E.1.42} cancels the boundary term in
Eq.~\eqref{eq:E.1.41}. Thus, when boundary terms are taken into account,
\(S'_G\) is the appropriate action to use for general relativity.
